\def\ProblemCode{SEGPERSLAZY}
\def\ProblemTitle{Lazy Persistent Segment Tree}
\makeproblem{\ProblemCode}{\ProblemTitle}

Cho dãy số nguyên $a_1, a_2, \dots, a_n$.
Ban đầu, tồn tại phiên bản $1$ của dãy, chính là dãy ban đầu.
Mỗi truy vấn cập nhật sẽ tạo ra một phiên bản mới của dãy.

Xử lý $q$ truy vấn thuộc một trong các dạng sau:

\begin{itemize}
    \item \boxed{1\ v\ l\ r\ k}:
    Từ phiên bản $v$, tạo ra một phiên bản mới bằng cách tăng mỗi phần tử
    trong đoạn $[l, r]$ lên $k$ đơn vị $(|k| \le 10^{9})$.

    \item \boxed{2\ v\ l\ r}:
    Tính tổng các phần tử trong đoạn $[l, r]$ của phiên bản $v$.

    \item \boxed{3\ v\ l\ r}:
    Tìm giá trị nhỏ nhất trong đoạn $[l, r]$ của phiên bản $v$.

    \item \boxed{4\ v\ l\ r}:
    Tìm giá trị lớn nhất trong đoạn $[l, r]$ của phiên bản $v$.
\end{itemize}

Các phiên bản được đánh số liên tiếp từ $1$ theo thứ tự xuất hiện
(các truy vấn loại \texttt{1}).

% ===== INPUT DESCRIPTION =====
\inputDesc{}

\begin{itemize}
    \item Dòng đầu tiên gồm hai số nguyên dương $n, q\ (1 \le n, q \le 10^{5})$.
    \item Dòng thứ hai gồm $n$ số nguyên $a_1, a_2, \dots, a_n\ (|a_i| \le 10^{9})$.
    \item $q$ dòng tiếp theo, mỗi dòng chứa một truy vấn như mô tả ở trên.
\end{itemize}

% ===== OUTPUT DESCRIPTION =====
\outputDesc{}

\begin{itemize}
    \item Với mỗi truy vấn loại 2, 3 hoặc 4, in ra một dòng chứa kết quả tương ứng.
\end{itemize}

\vspace{0.5em}

\ExampleBox{2}{
6 9\\
2 -5 4 1 7 -3\\
2 1 1 6\\
1 1 2 5 3\\
2 2 1 6\\
3 2 3 6\\
4 2 1 4\\
1 2 4 6 -2\\
2 3 4 6\\
3 1 1 6\\
4 3 1 6
}{
6\\
18\\
-3\\
7\\
-3\\
-5\\
7
}{}
